% Verbatim from /Users/alingalatan/Downloads/cubic_chaos_partial_results (2).zip:mom3_large_cumulant.tex (16 November 2026); GPT (a ChatGPT session of the user), relayed by the user.
% /workspace/scratch/3efc7fe714fa/mom3_large_cumulant.tex
\documentclass[11pt]{article}
\usepackage{amsmath,amssymb,amsthm,geometry,hyperref}
\geometry{margin=1in}
\newtheorem{theorem}{Theorem}
\newtheorem{lemma}{Lemma}
\newcommand{\E}{\mathbb E}
\newcommand{\He}{\operatorname{He}}
\newcommand{\sym}{\operatorname{sym}}
\title{An explicit large-fourth-cumulant bound in the third Gaussian chaos}
\author{}
\date{}
\begin{document}
\maketitle

All statements below are rigorous. The one-dimensional integral inequalities
in Lemma~\ref{lem:certificate} have a computer-assisted proof using only exact
integer arithmetic; the accompanying self-contained Python file is
\texttt{mom3\_interval\_certificate.py}. No floating-point numerical integration
is used in the proof. These results do not settle the intermediate-cumulant
regime of the full conjecture.

Write $c=\sqrt{2/\pi}$ and normalize $F\in\mathcal H_3$ by $\E F^2=1$.
Let $\delta=\E F^4-3$. Define, for $74\leq\delta\leq90$,
\[
 v(\delta)=\frac{3\sqrt{\delta/45-1}-1}{2},
 \qquad
 K(v)=
 \begin{cases}
 \displaystyle\frac{973}{1000}+\frac{29}{50}\left(\frac34-v\right),
       &v\leq\frac34,\\[4pt]
 \displaystyle\frac{973}{1000}-\frac12\left(v-\frac34\right),
       &v\geq\frac34.
 \end{cases}
\]

\begin{theorem}\label{thm:main}
For every dimension, every normalized $F\in\mathcal H_3$ with $\delta\geq74$
satisfies
\[
 \E|F|<cK(v(\delta))<\frac{3124}{3125}\,c<c.
\]
In particular, throughout the stronger regime $\delta\geq305/4$,
\[
 \E|F|<c\left[\frac{973}{1000}
       -\frac12\left(v(\delta)-\frac34\right)\right]
 \leq\frac{973}{1000}\,c.
\]
\end{theorem}

\paragraph{Tensor normalization and the fourth cumulant.}
Represent $F=I_3(f)$ for a real symmetric tensor $f$, with the normalization
$\E I_q(g)^2=q!\|g\|^2$. Thus $\|f\|^2=1/6$.
Contractions are taken in the underlying real Gaussian Hilbert space. Put
\[
 A=f\otimes_2f,\qquad a=\|A\|_{\mathrm{HS}}^2,
 \qquad b=\|\sym(f\otimes_1f)\|^2.
\]
Here $A$ is positive semidefinite and $\operatorname{tr}A=\|f\|^2=1/6$.
The product formula gives
\[
 F^2=I_6(\sym(f\otimes f))
       +9I_4(\sym(f\otimes_1f))+18I_2(f\otimes_2f)+6\|f\|^2.
\]
To evaluate the first term's variance, grouping permutations according to
the number of indices exchanged between the two triples gives
\[
 6!\|\sym(f\otimes f)\|^2
 =(3!)^2\sum_{r=0}^3\binom3r^2\|f\otimes_rf\|^2
 =72\|f\|^4+648a.
\]
We used $\|f\otimes_1f\|^2=\|f\otimes_2f\|^2$: the two contractions are
the Gram operators $T^*T$ and $TT^*$ for the flattening of $f$ into one
index versus two indices. Orthogonality of the chaos terms now gives the exact identity
\[
 \delta=1296a+1944b.
\]
Symmetrization is an orthogonal projection, so $b\leq a$. Consequently,
with $C=6A$,
\begin{equation}\label{eq:cumulant}
 C\succeq0,\qquad\operatorname{tr}C=1,
 \qquad \delta\leq90\operatorname{tr}(C^2)\leq90.
\end{equation}
All these formulas also hold for the Hilbert-space definition of a closed
Wiener chaos, by approximation, although the requested finite-dimensional
setting already suffices.

\paragraph{A large fourth cumulant forces a large one-coordinate cubic component.}
Let $p$ be the largest eigenvalue of $C$, and let $e$ be a corresponding unit
eigenvector. Positivity and trace one imply
\[
 \operatorname{tr}(C^2)\leq p,
 \qquad\operatorname{tr}(C^2)\leq p^2+(1-p)^2.
\]
If $\delta>45$, the first inequality gives $p>1/2$; combining the second
with \eqref{eq:cumulant} therefore yields
\begin{equation}\label{eq:p}
 p\geq\frac{1+\sqrt{\delta/45-1}}2.
\end{equation}
The normalized tensor $\psi=\sqrt6 f$ has one-coordinate marginal $C$.
Let $P$ be the orthogonal projection onto $\mathbb Re$. The three copies
of $P$ on the tensor factors commute, and
\[
 P^{\otimes3}\succeq I-\sum_{j=1}^3(I-P)_j.
\]
Taking its quadratic form at $\psi$ gives
\begin{equation}\label{eq:u}
 u:=6f(e,e,e)^2
   =\langle\psi,P^{\otimes3}\psi\rangle\geq3p-2\geq v(\delta).
\end{equation}

\paragraph{Conditional second moments improve the triangle bound.}
Write $X=G(e)$ and let $Y$ denote the independent Gaussian coordinates in
$e^\perp$. The exact chaos decomposition is
\[
 F=\alpha\He_3(X)+(X^2-1)L(Y)+XQ(Y)+R(Y),
\]
where $L\in\mathcal H_1(Y)$, $Q\in\mathcal H_2(Y)$, and
$R\in\mathcal H_3(Y)$. Set
\[
 u=6\alpha^2,\qquad\beta=2\E L^2,\qquad\gamma=\E Q^2,
 \qquad\eta=\E R^2.
\]
These quantities are nonnegative and $u+\beta+\gamma+\eta=1$.
Define the nonnegative polynomials
\[
 a_0(x)=\frac{(x^3-3x)^2}{6},\qquad
 b_0(x)=\frac{(x^2-1)^2}{2}.
\]
Orthogonality of the different $Y$-chaoses cancels every conditional cross
term. Conditional Cauchy--Schwarz therefore gives
\begin{equation}\label{eq:jensen}
 \E|F|\leq\E\sqrt{u a_0(X)+\beta b_0(X)+\gamma X^2+\eta}.
\end{equation}
No independence between $L,Q,R$ is asserted or needed.

\begin{lemma}[Exact integral certificate]\label{lem:certificate}
Let $D(x)=3a_0(x)/4+1/4$ and define four unnormalized half-Gaussian integrals
\begin{align*}
 J&=\int_0^\infty\sqrt{D(x)}e^{-x^2/2}\,dx,\\
 T_a&=\int_0^\infty\frac{a_0(x)-1}{\sqrt{D(x)}}e^{-x^2/2}\,dx,\\
 T_b&=\int_0^\infty\frac{b_0(x)-1}{\sqrt{D(x)}}e^{-x^2/2}\,dx,\\
 T_c&=\int_0^\infty\frac{x^2-1}{\sqrt{D(x)}}e^{-x^2/2}\,dx.
\end{align*}
Then
\[
 J<\frac{973}{1000},\qquad
 -\frac{29}{25}<T_a<-1,\qquad T_b<0,\qquad T_c<0.
\]
\end{lemma}

\begin{proof}[Exact arithmetic verification]
The accompanying script implements directed integer interval arithmetic at
scale $2^{96}$. On $[0,8]$ it partitions into intervals of width $1/2000$.
The cubic $x^3-3x$ is monotone on each cell, since its only nonnegative
critical point is the grid point $1$. Its square is enclosed including
zero whenever its range crosses zero. The same interval rules enclose
$(x^2-1)^2$, and $D\geq1/4$ ensures that every division is by a positive
interval.

For a density endpoint $e^{-t}$ with $0\leq t\leq32$, write
$e^{-t}=(e^{-t/64})^{64}$. The degree-15 alternating Taylor sum at
$t/64\leq1/2$ is a lower bound and the degree-16 sum an upper bound.
Six interval squarings preserve this enclosure. Integer square roots with
outward rounding enclose the other square roots. Each cell integral is
enclosed by the cell length times the entire integrand interval.

For the tails, if $x\geq8$ then $\sqrt D\leq x^3+1/2$, so its integral tail
is bounded by
\[
 \left(8^2+2+\frac1{16}\right)e^{-32}<\frac{67}{2^{32}}.
\]
The absolute value of each of the other three integrands, before multiplication
by the density, is at most $x^6$. Gaussian integration by parts gives
\[
 \int_8^\infty x^6e^{-x^2/2}\,dx
 \leq\left(8^5+5\cdot8^3+15\cdot8+\frac{15}{8}\right)e^{-32}
 <\frac{36000}{2^{32}}.
\]
Here only $e>2$ is used. The program sums the rigorous cell and tail intervals
and verifies the four rational inequalities in the statement by exact integer
comparisons. For readability, coarsened decimal enclosures implied by those
exact bounds are
\[
 \begin{array}{c|c}
 J &(0.9713232,\ 0.9722280)\\
 T_a&(-1.1514429,\ -1.1491778)\\
 T_b&(-0.6847749,\ -0.6817016)\\
 T_c&(-0.3355140,\ -0.3326555).
 \end{array}
\]
The decimal display is not used to perform any of the verification.
\end{proof}

\begin{proof}[Proof of Theorem~\ref{thm:main}]
The tangent inequality $\sqrt z\leq\sqrt d+(z-d)/(2\sqrt d)$ for $d>0$
applied in \eqref{eq:jensen} with $d=D(X)$ gives, since
$\eta=1-u-\beta-\gamma$,
\[
 \frac{\E|F|}{c}
 \leq J+\frac{u-3/4}{2}T_a+\frac\beta2T_b+\frac\gamma2T_c.
\]
The factor $c$ comes from integrating an even function against the standard
Gaussian density. Lemma~\ref{lem:certificate} makes the last two terms
nonpositive. Its upper and lower bounds on $T_a$ show that the right-hand
side is strictly less than $K(u)$. The function $K$ is continuous and
decreasing, and \eqref{eq:u} therefore gives
$\E|F|/c<K(v(\delta))$.

For $\delta\geq74$, one has
\[
 v(\delta)\geq\frac{3\sqrt{29/45}-1}{2}>\frac{88}{125},
\]
because $29/45-(301/375)^2=24/140625>0$.
Thus
\[
 K(v(\delta))<K(88/125)
 =\frac{973}{1000}+\frac{29}{50}\cdot\frac{23}{500}
 =\frac{3124}{3125}<1.
\]
Finally, $v(305/4)=3/4$, giving the second displayed conclusion.
\end{proof}

\paragraph{Why the first six moment bounds alone cannot prove the conjecture.}
This is a moment-relaxation obstruction, not a chaos counterexample.
Let a symmetric random variable $Z$ have absolute values
\[
 s_0=\frac9{10},\qquad t_0=\sqrt{\frac{9219}{19}},
 \qquad \mathbb P(|Z|=t_0)=\frac{361}{920361},
\]
with independent equiprobable signs. Direct algebra gives
\[
 \E Z^2=1,\qquad\E Z^4=93,\qquad
 \E Z^6=\frac{21283272}{475}<53745,
 \qquad\E|Z|>\frac9{10}>c.
\]
For comparison, the normalized $\He_3(G)/\sqrt6$ has fourth moment $93$
and sixth moment $53745$, as follows by expansion and Gaussian even moments.
Hence a proposed proof that uses only normalization, the fourth moment, and
the uniform sixth-moment upper bound would also apply to $Z$ and cannot
establish the desired sharp $L^1$ inequality. The tensor structure used
above supplies essential information beyond those scalar moment constraints.

\end{document}
